\documentclass[a4paper,12pt]{article}
\usepackage[italian]{babel}
\usepackage[utf8]{inputenc}
\usepackage[T1]{fontenc}
\usepackage{geometry}
\usepackage{hyperref}

\geometry{margin=2.5cm}

\title{\textbf{Design e Progettazione}\\
Elaborato Tecnologie Web 25/26}
\author{Casali Marco \\ Mancini Nicolas \\ Pulga Luca}
\date{22/01/2026}

\begin{document}

\maketitle

\begin{center}
\textbf{Unibo MatchSkills}
\end{center}

\newpage
\tableofcontents
\newpage


\section{Introduzione e problema}

Nel contesto universitario e accademico, la collaborazione tra studenti rappresenta un elemento fondamentale per il successo formativo e professionale. Tuttavia, la creazione di gruppi di studio efficaci o di team per progetti, competizioni e hackathon risulta spesso complessa. Le principali difficoltà risiedono nella ricerca di compagni con competenze adeguate, nella gestione delle diverse disponibilità temporali e nell’allineamento sugli obiettivi comuni.

Queste problematiche non si limitano al solo ambito universitario, ma si estendono anche ad altre istituzioni scolastiche secondarie, dove manca spesso uno strumento strutturato in grado di facilitare la collaborazione e la condivisione delle competenze.

\section{User Personas e Scenarios}

Per comprendere in modo approfondito le esigenze degli utenti e i diversi contesti di utilizzo della piattaforma, sono state definite specifiche User Personas e relativi Scenarios. Questo lavoro ha permesso di individuare i principali profili di studenti, le loro motivazioni e le modalità con cui interagiscono con il sistema.

La documentazione completa è disponibile all’interno del repository nel file:\\
\texttt{Unibo-Events-Management/docs/UniBo MatchSkills - User Personas e Scenarios.pdf}.

\section{Soluzione proposta}

La soluzione individuata è la realizzazione di una piattaforma web denominata \textbf{Unibo MatchSkills}. Il sistema è pensato principalmente per supportare gli studenti dell’Università di Bologna, ma è progettato in modo sufficientemente flessibile da poter essere adottato anche da altre istituzioni scolastiche.

L’obiettivo della piattaforma è facilitare la creazione di connessioni collaborative efficaci, favorendo la nascita di gruppi di studio, team di progetto e squadre per la partecipazione a competizioni e hackathon. Unibo MatchSkills nasce quindi dall’esigenza di superare le difficoltà tipiche della collaborazione accademica, offrendo un ambiente digitale centralizzato e strutturato.

In una futura evoluzione del sistema, non ancora disponibile nella prima release, è prevista l’integrazione di un sistema di matchmaking intelligente. Tale strumento sarà in grado di analizzare i profili degli studenti e suggerire connessioni ottimali sulla base di competenze tecniche, soft skill e disponibilità temporali.

Oltre alle funzionalità di matchmaking, la piattaforma consente agli utenti di creare e promuovere eventi interni, come coding challenge, workshop tra studenti o sessioni di studio collaborative. Inoltre, permette la creazione di team dedicati alla partecipazione a eventi sia interni che esterni e offre la possibilità di dichiarare l’intenzione di partecipare a un evento, facilitando la ricerca attiva di nuovi membri.

Unibo MatchSkills si configura quindi come un vero e proprio hub collaborativo, capace di incentivare la crescita delle competenze, la partecipazione attiva alla vita accademica e lo sviluppo di esperienze pratiche e competitive, fondamentali per il percorso formativo e professionale degli studenti.

\section{Casi d’Uso}

Le funzionalità della piattaforma sono state modellate attraverso una serie di casi d’uso, che descrivono le principali interazioni tra gli utenti e il sistema. I diagrammi dei casi d’uso sono disponibili online e all’interno del repository di progetto.

In particolare, è possibile consultare il diagramma completo al seguente link:\\
\url{https://lucid.app/lucidchart/b53577fd-765b-4d06-9503-2e2527c6091d}.

Il documento PDF è inoltre disponibile nel file:\\
\texttt{Unibo-Events-Management/docs/UniBo MatchSkills – Use Cases.pdf}.

\section{Design e Interfaccia Utente (UX/UI)}

La progettazione dell’interfaccia utente ha seguito principi di semplicità, coerenza visiva e accessibilità. Per quanto riguarda la scelta cromatica, è stata adottata una palette ispirata a Virtuale Unibo, in modo da mantenere una coerenza visiva con l’identità grafica istituzionale.

Il font selezionato è \textbf{Montserrat}, scelto per le sue caratteristiche di leggibilità e per il buon livello di accessibilità su diversi dispositivi. Le icone utilizzate provengono dal \textit{Material Design System}, garantendo uno stile moderno e riconoscibile.

Un aspetto centrale del processo di design è stato l’approccio \textit{mobile-first}. Per supportare il principio di \textit{progressive enhancement}, il design è stato inizialmente pensato per dispositivi mobili, utilizzando come riferimento uno smartphone con schermo di dimensioni ridotte, nello specifico l’iPhone SE. Successivamente, l’interfaccia è stata adattata a schermi di dimensioni maggiori.

I mockup dell’applicazione sono disponibili nel file:\\
\texttt{Unibo-Events-Management/docs/Unibo-MatchSkills.fig}.

\section{Progettazione del Database}

La progettazione del database è stata realizzata con l’obiettivo di supportare in modo efficiente la gestione degli utenti, dei ruoli, dei gruppi, degli eventi e delle competenze. La struttura dati è pensata per garantire flessibilità e scalabilità, permettendo future estensioni della piattaforma, in particolare in relazione al sistema di matchmaking.

\section{Funzionalità principali}

All’interno di Unibo MatchSkills, gli utenti assumono ruoli diversi in base alle attività svolte e al contesto di utilizzo. Lo studente rappresenta l’utente principale del sistema e può ricoprire più ruoli contemporaneamente o in momenti differenti della sua esperienza sulla piattaforma.

Quando uno studente crea un evento interno, fonda un gruppo di studio o avvia un team di progetto, assume il ruolo di \textit{Lead}. In questo caso è responsabile dell’organizzazione delle attività, della descrizione dell’evento o del gruppo e della gestione dei partecipanti. Al contrario, quando prende parte a eventi o team creati da altri utenti, lo studente assume il ruolo di \textit{Participant}.

Ogni studente dispone di un profilo personale strutturato e personalizzabile, che costituisce la base del sistema di matchmaking. Il profilo include informazioni relative alle competenze tecniche, come linguaggi di programmazione, framework e strumenti utilizzati, oltre alle soft skill, tra cui lavoro di squadra, comunicazione, leadership e problem solving.

La piattaforma consente inoltre la creazione di diversi tipi di gruppi collaborativi, come gruppi di studio, team per progetti accademici o team dedicati alla partecipazione a eventi esterni. Ogni gruppo è caratterizzato da un creatore, da un insieme di membri e da uno stato che indica se il gruppo è in fase di ricerca di nuovi partecipanti, completo o inattivo.

Unibo MatchSkills permette anche l’organizzazione e la promozione di eventi interni all’istituzione, quali mini hackathon, coding challenge, workshop tra pari e sessioni di studio collaborative. Ogni evento include una descrizione dettagliata, informazioni temporali e logistiche, eventuali competenze richieste e tag tematici per facilitarne la ricerca e la raccomandazione.

Per quanto riguarda gli eventi esterni, come hackathon e competizioni, la piattaforma consente agli studenti di cercare opportunità provenienti da piattaforme dedicate o inserite manualmente. Gli utenti possono dichiarare la propria intenzione di partecipare e utilizzare il sistema per trovare altri studenti con cui completare il team.

\section{Pannello admin}

È previsto un pannello di amministrazione dedicato agli utenti con ruolo di admin. Attraverso questo pannello è possibile moderare gruppi ed eventi, gestire il sistema di competenze e generare statistiche e report sull’utilizzo della piattaforma. Questi strumenti risultano fondamentali per monitorare l’andamento del sistema e supportarne il miglioramento continuo.

\section{Futura integrazione: matchmaking intelligente}

Il cuore evolutivo della piattaforma è rappresentato dal sistema di matchmaking intelligente. Tale sistema è progettato per suggerire connessioni rilevanti tra studenti, gruppi ed eventi, favorendo la creazione di collaborazioni efficaci.

Il meccanismo di raccomandazione potrà includere il calcolo di un punteggio di compatibilità basato su competenze, obiettivi formativi e disponibilità, oltre all’utilizzo di tecniche di intelligenza artificiale e all’analisi delle interazioni precedenti degli utenti sulla piattaforma.

\end{document}
